% Part: first-order-logic
% Chapter: completeness
% Section: maximally-consistent-sets

% Definition of maximally consistent sets. Properties of mcs required
% for completeness proved are in maximally-consistent-sets.tex in the
% chapter on the proof system used.

\documentclass[../../../include/open-logic-section]{subfiles}

\begin{document}

\olfileid{fol}{com}{mcs}
\olsection{Himpunan \usetoken{P}{sentence} kang Konsisten Maksimal}

\begin{defn}[Himpunan konsisten maksimal]
\ollabel{mcs}
Himpunan~$\Gamma$ saka !!{sentence}s iku
\emph{konsisten maksimal} yèn lan mung yèn
\begin{enumerate}
\item $\Gamma$ konsisten, lan
\item yèn $\Gamma \subsetneq \Gamma'$, mula $\Gamma'$ ora konsisten.
\end{enumerate}
\end{defn}

\begin{explain}
Definisi liya kang padha tegesé yaiku: himpunan $\Gamma$
saka ukara iku \emph{konsisten maksimal} yèn lan mung yèn
\begin{enumerate}
\item $\Gamma$ konsisten, lan
\item yèn $\Gamma \cup \{ !A \}$ konsisten, mula $!A \in \Gamma$.
\end{enumerate}
Kanthi tembung liya, ora bisa nambah !!{sentence}s kang
durung ana ing~$\Gamma$ marang himpunan konsisten maksimal~$\Gamma$
tanpa njalari himpunan kang luwih gedhé dadi ora konsisten.
\end{explain}

\begin{explain}
Himpunan konsisten maksimal wigati ing bukti kelengkapan,
amarga saben himpunan konsisten saka !!{sentence}s~$\Gamma$
bisa dilebokaké ing sawijining himpunan konsisten maksimal~$\Gamma^*$.
Kajaba iku, kanggo saben !!{sentence}~$!A$, himpunan konsisten
maksimal ngemot $!A$ utawa negasiné $\lnot !A$.
Iki uga bener kanggo !!{sentence}s atomik. Mula, saka himpunan
konsisten maksimal ing basa kang wis ditambahi !!{constant}s
kanthi cocog, kita bisa mbangun !!a{structure} kang interpretasi
!!{predicate}s-é ditemtokaké déning !!{sentence}s atomik endi
kang ana ing~$\Gamma^*$. Banjur bisa dituduhaké yèn !!{structure}
iki ndadèkaké kabèh !!{sentence}s ing~$\Gamma^*$
(lan mula uga ing $\Gamma$) bener. Bukti bab iki mbutuhaké
$\lnot !A \in \Gamma^*$ yèn lan mung yèn $!A \notin \Gamma^*$,
$(!A \lor !B) \in \Gamma^*$ yèn lan mung yèn $!A
\in \Gamma^*$ utawa $!B \in \Gamma^*$, lan sapituruté.
\end{explain}

\begin{prop}
\ollabel{prop:mcs}
Anggepen $\Gamma$ konsisten maksimal. Mula:
\begin{enumerate}
\item \ollabel{prop:mcs-prov-in} Yèn $\Gamma \Proves !A$, mula $!A \in
  \Gamma$.

\item \ollabel{prop:mcs-either-or} Kanggo saben $!A$, $!A \in
  \Gamma$ utawa $\lnot !A \in \Gamma$.

\tagitem{prvAnd}{\ollabel{prop:mcs-and} $(!A \land !B) \in \Gamma$
  yèn lan mung yèn $!A \in \Gamma$ lan $!B \in \Gamma$.}{}

\tagitem{prvOr}{\ollabel{prop:mcs-or} $(!A \lor !B) \in \Gamma$ yèn lan
  mung yèn $!A \in \Gamma$ utawa $!B \in \Gamma$.}{}

\tagitem{prvIf}{\ollabel{prop:mcs-if} $(!A \lif !B) \in \Gamma$ yèn lan
  mung yèn $!A \notin \Gamma$ utawa $!B \in \Gamma$.}{}
\end{enumerate}
\end{prop}

\begin{proof}
Kanggo sakabèhé ing ngisor iki, anggepen $\Gamma$
konsisten maksimal.
\begin{enumerate}
\item Yèn $\Gamma \Proves !A$, mula $!A \in \Gamma$.

Anggepen $\Gamma \Proves !A$. Kanggo mbantah, anggepen $!A
\notin \Gamma$. Amarga $\Gamma$ konsisten maksimal, $\Gamma
\cup \{!A\}$ ora konsisten, mula $\Gamma \cup \{!A\} \Proves
\lfalse$. Miturut
\tagrefs{prfSC/{fol:seq:prv:prop:provability-contr},prfND/{fol:ntd:prv:prop:provability-contr}},
$\Gamma$ ora konsisten. Iki mbantah anggepan yèn
$\Gamma$ konsisten. Mula, ora mungkin $!A \notin
\Gamma$, dadi $!A \in \Gamma$.

\item Kanggo saben $!A$, $!A \in \Gamma$ utawa $\lnot !A \in \Gamma$.

Kanggo mbantah, anggepen ana~$!A$ kang nyukupi $!A \notin \Gamma$
lan $\lnot !A \notin \Gamma$. Amarga $\Gamma$ konsisten maksimal,
$\Gamma \cup \{!A\}$ lan $\Gamma \cup \{\lnot !A\}$ padha-padha
ora konsisten, mula $\Gamma \cup \{!A\} \Proves \lfalse$ lan $\Gamma \cup
\{\lnot !A\} \Proves \lfalse$. Miturut
\tagrefs{prfSC/{fol:seq:prv:prop:provability-exhaustive},prfND/{fol:ntd:prv:prop:provability-exhaustive}},
$\Gamma$ ora konsisten, sawijining kontradiksi. Mula ora ana
!!a{sentence}~$!A$ kaya mangkono, lan kanggo saben $!A$,
$!A \in \Gamma$ utawa $\lnot !A \in \Gamma$.

\tagitem{defAnd}{}{%
\iftag{probAnd}{Gladhen.}{%
$(!A \land !B) \in \Gamma$ yèn lan mung yèn
$!A \in \Gamma$ lan $!B \in \Gamma$:

Kanggo arah maju, anggepen $(!A \land !B) \in \Gamma$.
Mula $\Gamma \Proves !A \land !B$. Miturut
\tagrefs{prfSC/{fol:seq:prv:prop:provability-land-left},prfND/{fol:ntd:prv:prop:provability-land-left}},
$\Gamma \Proves !A$ lan $\Gamma \Proves !B$. Miturut
\olref{prop:mcs-prov-in}, $!A \in \Gamma$ lan $!B \in \Gamma$,
kaya kang dikarepaké.

Kanggo arah walikan, anggepen $!A \in \Gamma$ lan $!B \in
\Gamma$. Mula $\Gamma \Proves !A$ lan $\Gamma \Proves !B$. Miturut
\tagrefs{prfSC/{fol:seq:prv:prop:provability-land-right},prfND/{fol:ntd:prv:prop:provability-land-right}},
$\Gamma \Proves !A \land !B$. Miturut \olref{prop:mcs-prov-in},
$(!A \land !B) \in \Gamma$.}}

\tagitem{defOr}{}{%
\iftag{probOr}{Gladhen.}{%
$(!A \lor !B) \in \Gamma$ yèn lan mung yèn
$!A \in \Gamma$ utawa $!B \in \Gamma$.

Kanggo kontrapositif arah maju, anggepen $!A
\notin \Gamma$ lan $!B \notin \Gamma$. Kita arep nuduhaké
$(!A \lor !B) \notin \Gamma$. Amarga $\Gamma$ konsisten maksimal,
$\Gamma \cup \{!A\} \Proves \lfalse$ lan $\Gamma \cup \{!B\} \Proves \lfalse$.
Miturut
\tagrefs{prfSC/{fol:seq:prv:prop:provability-lor-left},prfND/{fol:ntd:prv:prop:provability-lor-left}},
$\Gamma \cup \{(!A \lor !B)\}$ ora konsisten. Mula $(!A \lor !B)
\notin \Gamma$, kaya kang dikarepaké.

Kanggo arah walikan, anggepen $!A \in \Gamma$ utawa $!B \in
\Gamma$. Mula $\Gamma \Proves !A$ utawa $\Gamma \Proves !B$. Miturut
\tagrefs{prfSC/{fol:seq:prv:prop:provability-lor-right},prfND/{fol:ntd:prv:prop:provability-lor-right}},
$\Gamma \Proves !A \lor !B$. Miturut \olref{prop:mcs-prov-in},
$(!A \lor !B) \in \Gamma$, kaya kang dikarepaké.}}

\tagitem{defIf}{}{%
\iftag{probIf}{Gladhen.}{%
$(!A \lif !B) \in \Gamma$ yèn lan mung yèn
$!A \notin \Gamma$ utawa $!B \in \Gamma$:

Kanggo arah maju, anggepen $(!A \lif !B) \in \Gamma$,
lan kanggo mbantah anggepen $!A \in \Gamma$ sarta
$!B \notin \Gamma$. Saka anggepan iki,
$\Gamma \Proves !A \lif !B$ lan $\Gamma \Proves !A$. Miturut
\tagrefs{prfSC/{fol:seq:prv:prop:provability-mp},prfND/{fol:ntd:prv:prop:provability-mp}},
$\Gamma \Proves !B$. Nanging miturut \olref{prop:mcs-prov-in},
$!B \in \Gamma$, mbantah anggepan yèn $!B \notin \Gamma$.

Kanggo arah walikan, dhisik delengen kahanan $!A \notin
\Gamma$. Miturut \olref{prop:mcs-either-or},
$\lnot !A \in \Gamma$, mula $\Gamma \Proves \lnot !A$. Miturut
\tagrefs{prfSC/{fol:seq:prv:prop:provability-lif},prfND/{fol:ntd:prv:prop:provability-lif}},
$\Gamma \Proves !A \lif !B$. Miturut \olref{prop:mcs-prov-in}
sapisan manèh, kita olèh $(!A \lif !B) \in \Gamma$,
kaya kang dikarepaké.

Saiki delengen kahanan $!B \in \Gamma$. Mula $\Gamma \Proves !B$,
lan miturut
\tagrefs{prfSC/{fol:seq:prv:prop:provability-lif},prfND/{fol:ntd:prv:prop:provability-lif}},
$\Gamma \Proves !A \lif !B$. Miturut \olref{prop:mcs-prov-in},
$(!A \lif !B) \in \Gamma$.}}
\end{enumerate}
\end{proof}

\begin{prob}
Rampungna bukti kanggo \olref[fol][com][mcs]{prop:mcs}.
\end{prob}

\end{document}
